% 钱德拉塞卡拉·拉曼（综述）
% license CCBYSA3
% type Wiki

本文根据 CC-BY-SA 协议转载翻译自维基百科\href{https://en.wikipedia.org/wiki/C._V._Raman}{相关文章}

钱德拉塞卡拉·文卡塔·“C. V.”·拉曼爵士（Sir Chandrasekhara Venkata "C. V." Raman，/ˈrɑːmən/，\(^\text{[2]}\)音译“拉曼”；泰米尔语：சந்திரசேகர வெங்கட ராமன்，罗马化：Cantiracēkara Veṅkaṭa Rāmaṉ；1888年11月7日—1970年11月21日）是一位印度物理学家\(^\text{[3]}\)，以其在光散射领域的研究而闻名。\(^\text{[4]}\)他使用自己研制的分光仪，与学生K. S. 克里希南共同发现：当光线穿过透明物质时，被偏转的光会改变其波长。这种此前未知的光散射现象，他们称之为“修饰散射”，后来被命名为“拉曼效应”或“拉曼散射”。1930年，拉曼因这一发现获得诺贝尔物理学奖，成为第一位获得该奖项的亚洲人和非白人。\(^\text{[5]}\)

拉曼出生于泰米尔婆罗门家庭，自幼聪颖过人。他分别在11岁和13岁时，就完成了圣阿洛伊修斯英印高级中学的中学和高中课程。16岁时，他以优异成绩从马德拉斯大学附属的总督学院毕业，并以物理学荣誉学士的身份名列榜首。他的第一篇研究论文《光的衍射》发表于1906年，当时他仍是一名研究生。次年，他获得了硕士学位。19岁时，拉曼进入加尔各答的印度财政服务部门，担任副会计长。在那里，他结识了印度科学促进会（Indian Association for the Cultivation of Science，简称 IACS）——印度首个科研机构。该机构允许他独立开展研究，也正是在这里，拉曼在声学与光学领域作出了重要贡献。

1917年，拉曼受阿修托什·穆克吉任命，出任加尔各答大学下属拉贾巴扎尔理学院的首任帕利特物理学教授。在他首次前往欧洲的旅途中，亲眼目睹地中海的蓝色使他产生了强烈的好奇心——他认为当时流行的解释（即海洋的蓝色是由于天空中瑞利散射光的反射所致）是错误的。1926年，他创办了《印度物理学杂志》。1933年，他迁往班加罗尔，成为印度科学研究所的首位印度籍所长。同年，他还创立了印度科学院。1948年，他建立了拉曼研究所，并在此工作直至生命的最后时光。

拉曼效应于1928年2月28日被发现。印度政府将这一天定为“全国科学日”，每年举行纪念活动。
\subsection{早年生活与教育}
拉曼出生于英属印度马德拉斯省的蒂鲁吉拉帕利（Tiruchirappalli，今印度泰米尔纳德邦的蒂鲁吉拉帕利），父母均为泰米尔·艾耶尔婆罗门——父亲钱德拉塞卡拉·拉曼纳森·艾耶，\(^\text{[6][7]}\)母亲帕尔瓦蒂·阿玛尔。\(^\text{[8]}\)他是家中八个孩子中的第二个。\(^\text{[9]}\)他的父亲是一名中学教师，收入微薄。拉曼曾幽默地回忆说：“我出生时嘴里含着一把铜勺。因为在我出生时，父亲那‘辉煌’的薪水是每月十卢比！”\(^\text{[10]}\)1892年，全家迁往安得拉邦的维萨卡帕特南（当时称为Vizagapatam或Vizag），因为他的父亲被任命为A.V. 纳拉辛哈·饶夫人学院物理系教师。\(^\text{[11]}\)

拉曼在维萨卡帕特南的圣阿洛伊修斯英印高级中学接受教育。\(^\text{[12]}\)他在11岁时通过了中学毕业考试，13岁时以奖学金身份通过了艺术初级考试（First Examination in Arts，相当于今日的大学预科课程），\(^\text{[9][13]}\)并在两项考试中均名列安得拉邦学校委员会（今安得拉邦中等教育委员会）榜首。\(^\text{[14]}\)

1902年，拉曼进入马德拉斯（今钦奈）的总督学院就读，当时他的父亲已调任至该校教授数学与物理。\(^\text{[15]}\)1904年，他以优异成绩从马德拉斯大学获得文学学士（B.A.）学位，并在物理学与英语两门科目中均名列第一，荣获金牌奖。\(^\text{[14]}\)18岁时，仍为研究生的拉曼在英国期刊《哲学杂志》上发表了他的第一篇科学论文——《由矩形孔引起的不对称衍射带》（Unsymmetrical diffraction bands due to a rectangular aperture，1906年）。\(^\text{[16]}\)1907年，他以最高荣誉从同一所大学获得文学硕士（M.A.）学位。\(^\text{[17][18]}\)同年，他的第二篇论文发表于同一杂志，题为《液体的表面张力》。\(^\text{[19]}\)巧合的是，这篇论文与瑞利勋爵的一篇关于人耳对声音敏感度的论文同时刊登。\(^\text{[20]}\)由此，瑞利勋爵开始与拉曼通信，并礼貌地称呼他为“教授”。\(^\text{[14]}\)

拉曼的物理导师理查德·卢埃林·琼斯深知他的学术潜力，坚持建议他前往英国继续研究。\(^\text{[21]}\)琼斯为他安排了与吉法德上校的体检。然而，拉曼的健康状况欠佳，被认为是“体弱多病之人”。\(^\text{[22]}\)体检结果显示他无法适应英国严酷的气候。\(^\text{[11]}\)对此，拉曼后来回忆道：“吉法德检查了我，并证明如果我去英国，我会死于肺结核。”\(^\text{[23]}\)
\subsection{职业生涯}
拉曼的哥哥钱德拉塞卡拉·苏布拉马尼亚·艾耶早年已加入享有盛誉的印度政府部门——印度财政服务处（Indian Finance Service，今为印度审计与会计服务 Indian Audit and Accounts Service）。\(^\text{[24]}\)拉曼也步其后尘，于1907年2月的入职考试中名列第一，顺利取得任职资格。\(^\text{[25]}\)同年6月，他被派往加尔各答（今加尔各答）任副会计长。\(^\text{[11]}\)

在加尔各答期间，拉曼对印度科学促进会（Indian Association for the Cultivation of Science，简称 IACS）深感钦佩。\(^\text{[23]}\)该机构创立于1876年，是印度历史上第一所科学研究机构。他很快结识了几位重要人物：阿修托什·德伊，后来的长期合作伙伴；阿姆里塔·拉尔·瑟卡尔），IACS创办人马亨德拉拉尔·瑟卡尔之子、协会秘书；阿修托什·穆克吉，IACS执行委员、加尔各答大学校长。\(^\text{[14]}\)在他们的支持下，拉曼获准利用业余时间在IACS进行科学研究——哪怕是在“极其不寻常的时间”，\(^\text{[26]}\)正如他后来回忆的那样。\(^\text{[11]}\)当时，IACS尚未正式招募全职研究人员，也从未发表过科研论文。拉曼在1907年发表的论文《偏振光下的牛顿环》发表于《自然》杂志，成为该研究所的第一篇学术论文。\(^\text{[27]}\)这项成果极大地激发了IACS的科研热情，并促使其于1909年创办了《印度科学促进会公报》，其中拉曼是主要撰稿人。\(^\text{[14]}\)

1909年，拉曼被调任至英属缅甸仰光（今缅甸仰光），担任货币官。然而不久之后，他因父亲病逝而返回马德拉斯。父亲的去世及随后的丧礼使他在当地滞留了一整年。\(^\text{[28]}\)不久，他重新回到仰光任职，但在1910年又被调回印度马哈拉施特拉邦的那格浦尔。\(^\text{[29]}\)尚未在那格浦尔任职满一年，1911年他便被晋升为会计长，再次调往加尔各答。\(^\text{[28]}\)

从1915年起，加尔各答大学开始将研究生派至印度科学促进会（IACS），在拉曼的指导下进行科研工作。他的第一位学生是苏丹舒·库马尔·班纳吉（Sudhangsu Kumar Banerji），当时是加内什·普拉萨德指导下的博士研究生，后来成为印度气象局天文台总监。

自次年起，其他大学也纷纷效仿，如阿拉哈巴德大学、仰光大学、印多尔皇后学院、那格浦尔科学研究所、克里什纳斯学院以及马德拉斯大学等，\(^\text{[30]}\)都派学生前往IACS接受拉曼的指导。到1919年为止，拉曼已指导过十多名学生。\(^\text{[31]}\)同年，IACS秘书阿姆里塔·拉尔·瑟卡尔去世后，拉曼被授予两个荣誉职务：名誉教授与名誉秘书。\(^\text{[26]}\)拉曼将这一时期称为他人生的“黄金时代”。\(^\text{[32]}\)

加尔各答大学于1913年选定拉曼担任“帕利特物理学教授”，该职位由捐助人塔拉克纳特·帕利特爵士）设立。大学参议院于1914年1月30日正式任命，如会议记录所载：\\\\
“在1914年1月30日的参议院会议上，作出以下帕利特教授职位的任命：P. C. 雷博士与C. V. 拉曼先生。……每位教授的任命为长期制。教授须于年满六十岁时离任。”\(^\text{[14]}\)

在1914年之前，阿修托什·穆克吉曾邀请贾格迪什·钱德拉·玻色出任该职位，但玻色婉拒了邀请。\(^\text{[14]}\)作为第二人选，拉曼被任命为首任“帕利特物理学教授”，但由于第一次世界大战的爆发，他的上任被迫推迟。直到1917年，他加入加尔各答大学于1914年新设的拉贾巴扎尔理学院，才正式成为全职教授。\(^\text{[14]}\)拉曼在担任政府公务员十年后，勉强辞去职务——这一决定被称为一种“至高的牺牲”\(^\text{[26]}\)，因为教授的薪水几乎只有他原先收入的一半。不过，他的任职条件在加入大学时被明确写入报告，对他极为有利。报告中指出：\\\\
“C. V. 拉曼先生接受了塔拉克纳特·帕利特爵士物理学教授一职，条件是他无需被派往印度以外的地区……报告显示，C. V. 拉曼先生自1917年7月2日起正式就任帕利特物理学教授……拉曼先生已通知，他不必承担硕士或理学硕士课程的教学任务，以免影响其科研工作或指导高年级研究生进行研究。”\(^\text{[30]}\)

拉曼被任命为“帕利特物理学教授”一事，曾遭到加尔各答大学参议院部分成员的强烈反对，尤其是外国成员——理由是他既没有博士学位，也从未出国留学。作为回应，阿修托什·穆克吉特地安排加尔各答大学于1921年授予拉曼名誉理学博士学位，以示肯定。同年，他受邀前往牛津大学，在“大英帝国大学大会”上发表演讲。\(^\text{[34]}\)到那时，拉曼已在国际上享有盛誉，其接待方正是两位诺贝尔奖得主：J. J. 汤姆孙爵士与卢瑟福勋爵。\(^\text{[35]}\)1924年，拉曼当选为英国皇家学会院士后，穆克吉问他今后的目标，他笑答道：“当然是诺贝尔奖。”\(^\text{[26]}\)1926年，拉曼创办了《印度物理学杂志》，并担任首任主编。\(^\text{[36]}\)该杂志第二卷刊登了他著名的论文《一种新的辐射》，首次报道了“拉曼效应”的发现。\(^\text{[37][38]}\)

1932年，拉曼的“帕利特物理学教授”职位由德本德拉·莫汉·玻色接任。翌年，他被任命为位于班加罗尔的印度科学研究所（Indian Institute of Science，简称 IISc）所长，于1933年离开加尔各答。\(^\text{[39]}\)印度科学研究所的创立得到了多方支持：迈索尔国王克里希纳拉贾·瓦迪亚尔四世、实业家贾姆谢特吉·塔塔以及海得拉巴的尼扎姆——米尔·奥斯曼·阿里·汗爵士共同捐赠了土地与资金。1909年，印度总督明托勋爵批准建立该研究所，并由英国政府任命莫里斯·特拉弗斯为首任所长。\(^\text{[40]}\)拉曼成为第四任所长，也是第一位印度籍所长。在任期间，他聘用了G. N. 拉马钱德兰，后者后来成为著名的X射线晶体学家。1934年，拉曼创立了印度科学院，并创办了其学术期刊《印度科学院院刊》，该期刊后来分为三个分支刊物：《数学科学院刊》，《化学学报》，《地球系统科学学报》。\(^\text{[35]}\)大约在同一时期，拉曼早在1917年就提出设想的“加尔各答物理学会”也正式成立。\(^\text{[14]}\)

1943年，拉曼与他的前学生潘查帕凯萨·克里希纳穆尔蒂共同创办了“特拉凡哥尔化学制造有限公司”。\(^\text{[41]}\)该公司于1996年更名为“TCM有限公司”，是印度最早的有机与无机化学品制造企业之一。\(^\text{[42]}\)1947年，印度独立后新政府任命拉曼为首任“国家教授”。\(^\text{[43]}\)

拉曼于1948年从印度科学研究所退休，次年在班加罗尔创立了“拉曼研究所”。他担任所长，并一直在此工作，直至1970年去世。\(^\text{[43]}\)
\subsection{科学贡献}
\subsubsection{音乐声学}
\begin{figure}[ht]
\centering
\includegraphics[width=8cm]{./figures/3a7f4c6530b8aaba.png}
\caption{能级示意图：展示了参与拉曼信号的各个能级状态} \label{fig_QDLlm_1}
\end{figure}
拉曼的研究兴趣之一是音乐声音的科学基础。他受到赫尔曼·冯·亥姆霍兹所著《音的感觉》一书的启发，这本书是他加入印度科学促进会（IACS）后接触到的。\(^\text{[25]}\)在1916年至1921年间，拉曼大量发表了相关研究成果。他提出了基于速度叠加原理的弓弦乐器横向振动理论。他最早的研究之一是关于小提琴和大提琴中的“狼音”现象。\(^\text{[44][45]}\)此外，他还研究了各种小提琴及类似乐器的声学特性，包括印度传统弦乐器的声学行为。\(^\text{[46][47]}\)他甚至研究了水滴溅起的声音。\(^\text{[48]}\)拉曼还进行了他称之为“机械演奏小提琴实验”的研究。\(^\text{[49]}\)
\begin{figure}[ht]
\centering
\includegraphics[width=8cm]{./figures/8cb85b7f2b42ce80.png}
\caption{1930年诺贝尔奖颁奖典礼上的拉曼（左起第一位），与其他获奖者合影：C. V. 拉曼（物理学奖）、汉斯·费舍尔（化学奖）、卡尔·兰德施泰纳（医学奖）以及辛克莱·刘易斯（文学奖）。} \label{fig_QDLlm_2}
\end{figure}
拉曼还研究了印度鼓类乐器的独特性。\(^\text{[50]}\)他对塔布拉鼓与姆里当加姆鼓声音谐波特性的分析，是印度打击乐器声学的首批科学研究。\(^\text{[51]}\)他还撰写了一篇关于钢琴琴弦振动的批判性论文，探讨并修正了被称为考夫曼理论的观点。\(^\text{[52]}\)1921年短暂访问英国期间，拉曼抓住机会研究了伦敦圣保罗大教堂圆顶“回音廊”中的声音传播现象，该建筑能产生奇特的回声与声学效果。\(^\text{[53][54]}\)拉曼在声学方面的研究，无论在实验方法还是概念上，都为他后来的光学与量子力学研究奠定了重要基础。\(^\text{[55]}\)
\subsubsection{海水的蓝色}
在拓展光学研究的过程中，拉曼自1919年起开始研究光的散射现象。\(^\text{[56]}\)他在光学物理方面的第一个重要发现，便是对海水蓝色成因的解释。1921年9月，他从英国乘坐S.S. Narkunda号轮船返回印度途中，凝视地中海时，对海水的蓝色产生了浓厚兴趣。凭借随身携带的简单光学仪器——一台袖珍分光镜和一个尼科尔棱镜，他亲自对海水进行了观察与研究。\(^\text{[57]}\)当时，关于海水为何呈蓝色的多种假说中，\(^\text{[58][59]}\)最广为接受的是瑞利勋爵在1910年提出的解释：“深海那令人赞叹的深蓝色，与水本身的颜色无关，而只是天空蓝光在海面的反射。”\(^\text{[60]}\)瑞利对天空为何呈蓝色的解释是正确的——这一现象即如今所称的“瑞利散射”，即光在大气中被微粒散射与折射所致。\(^\text{[61][62]}\)因此，他关于海水蓝色来源于天空反射的观点在当时被普遍认为是正确的。然而，拉曼通过使用尼科尔棱镜消除了海面反射阳光的影响，直接观察海水后发现：海水在这种条件下显得比平时更加湛蓝。这一结果与瑞利的理论相矛盾，也使他确信，海水的蓝色并非仅仅来自天空的反射。\(^\text{[63]}\)

当S.S. Narkunda号轮船抵达孟买港（今孟买港）后，拉曼立即完成了一篇题为《海的颜色》的文章，并发表在1921年11月的《自然》杂志上。他在文中指出，瑞利的解释“通过一种简单的观察方式（使用尼科尔棱镜）就可以被质疑”。\(^\text{[63]}\)他写道：\\\\
“当通过一个放置在眼前的尼科尔棱镜观察海水，以消除表面反射时，可以看到阳光的光线穿入水中，并由于透视效应似乎在水的深处汇聚成一点。问题是：究竟是什么在散射光线，使得它的路径变得可见？一个值得考虑的有趣可能性是——这些散射光的粒子，至少部分，可能正是水分子本身。”\(^\text{[14]}\)
\begin{figure}[ht]
\centering
\includegraphics[width=6cm]{./figures/c373caa8734590ef.png}
\caption{拉曼《光的分子衍射》（Molecular Diffraction of Light，1922年）标题页} \label{fig_QDLlm_3}
\end{figure}
当拉曼抵达加尔各答后，他请来自仰光大学的学生K. R. 拉曼纳森在印度科学促进会（IACS）继续展开进一步研究。\(^\text{[64]}\)到1922年初，拉曼得出了自己的结论，并在《伦敦皇家学会学报》上发表了报告：\\\\
“本文拟提出一种完全不同的观点：在这一现象中，正如在天空呈色这一类似情形中一样，分子衍射决定了所观察到的亮度，并在很大程度上决定了其颜色。作为讨论的必要前提，本文将给出关于水中分子散射强度的理论计算与实验观察。”\(^\text{[65]}\)

正如拉曼所言，拉曼纳森于1923年发表了一篇详细的实验研究报告。\(^\text{[66]}\)随后，他于1924年对孟加拉湾的进一步研究提供了确凿的证据。\(^\text{[67]}\)如今人们已经知道，水本身的颜色主要源于其对光谱中红色与橙色长波光的选择性吸收——这种吸收来自水分子中O–H（氧与氢结合）伸缩振动模式的红外吸收泛频。\(^\text{[68]}\)
\subsubsection{拉曼效应}
\textbf{背景}

拉曼在光散射研究中的第二项重要发现，是一种全新的辐射现象——后来以他名字命名的“拉曼效应”。【69】在揭示了导致海水呈蓝色的光散射性质之后，拉曼将研究重点转向这一现象背后的原理。他在1923年的实验中发现：当阳光通过紫色滤光玻璃照射到某些液体和固体上时，除了入射光外，还会出现其他光线。这提示了可能存在新的散射形式。他的学生拉曼纳森当时认为这只是“荧光的微弱痕迹”。【14】1925年，新加入的研究助理K. S. 克里希南进一步指出，当光在液体中发生散射时，除了常见的偏振弹性散射外，还可能存在额外的一条散射谱线。\(^\text{[70]}\)他将这一现象称为“微弱的荧光”。\(^\text{[71]}\)然而，在接下来的两年中，对这一现象进行理论解释的尝试都未能取得成功。\(^\text{[72]}\)

这一研究的重大推动力来自“康普顿效应”的发现。1923年，圣路易斯华盛顿大学的阿瑟·康普顿发现了电磁波也可以用粒子来描述的证据。\(^\text{[73]}\)到1927年，这一现象已被科学界广泛接受，拉曼也认同这一观点。\(^\text{[74]}\)当年12月，康普顿获得诺贝尔物理学奖的消息传来时，拉曼兴奋地对克里希南说道：\\\\
“好消息……真是太好了！不过你看，克里希南——如果这对X射线成立，那么对可见光也一定成立。我一直这样认为。一定存在一个与康普顿效应对应的光学现象。我们必须去追寻它，而且我们正走在正确的道路上。它一定会被发现，而我们一定要拿下诺贝尔奖！”\(^\text{[75]}\)

然而，拉曼灵感的起源可以追溯得更远。正如康普顿后来回忆的那样：“也许正是那场多伦多的辩论，使他（拉曼）在两年后发现了拉曼效应。”\(^\text{[25]}\)所谓“多伦多辩论”，指的是1924年在多伦多举行的英国科学促进会会议上，关于“光量子是否存在”的讨论。会上，康普顿展示了自己的实验结果，而哈佛大学的威廉·杜安则提出相反的证据，主张光是波动。\(^\text{[76]}\)拉曼当时站在杜安一方，并直言道：“康普顿，你是个出色的辩手，但真理并不在你那一边。”\(^\text{[25]}\)

\textbf{散射实验}
\begin{figure}[ht]
\centering
\includegraphics[width=6cm]{./figures/ebf7f67e43fc5253.png}
\caption{拉曼与克里希南发表的早期苯的拉曼光谱\(^\text{[77]}\)} \label{fig_QDLlm_4}
\end{figure}
克里希南于1928年1月初开始了实验。\(^\text{[64]}\)1月7日，他发现无论使用哪种纯净液体，都会在可见光范围内产生偏振荧光。当拉曼看到这一结果时，他十分震惊——自己多年来竟从未观察到这种现象。\(^\text{[64]}\)当晚，他与克里希南将这一新现象命名为“修饰散射”，以区别于康普顿效应中的“未修饰散射”。2月16日，他们向《自然》杂志投稿了一篇题为《一种新的次级辐射》的论文，该文于3月31日发表。\(^\text{[78]}\)

1928年2月28日，他们首次成功获得了与入射光分离的“修饰散射”光谱。由于当时测定光波长的技术困难，他们最初依赖肉眼通过棱镜观察阳光所形成的颜色。拉曼为此发明了一种用于探测和测量电磁波的分光仪。\(^\text{[35][79]}\)谈到这项发明时，拉曼后来曾风趣地说：“当我获得诺贝尔奖时，我在设备上花的钱还不到200卢比。”\(^\text{[80]}\)不过，显然他在整个实验上的总花费远不止此数。\(^\text{[81]}\)从那一刻起，他们使用来自汞弧灯的单色光照射透明材料，并让透射光投射到分光仪上记录光谱。这样，散射光谱线便能够被精确测量与拍摄下来。\(^\text{[82][83]}\)

\textbf{公告}

同一天，拉曼向媒体正式宣布了这一发现。印度联合通讯社于次日，即2月29日报道了此事，标题为《辐射新理论：拉曼教授的发现》。\(^\text{[84]}\)报道写道：\\\\

“加尔各答大学的C. V. 拉曼教授作出了一项有望对物理学产生根本性意义的发现……
这一新现象所展现的特征，比康普顿教授用X射线所发现的还要令人震惊。其主要特征是：当物质被某一颜色的光激发时，其中的原子会发射出两种颜色的光，其中一种与激发光不同，且在光谱中位置更低。最令人惊讶的是，这种改变的颜色与所使用物质的种类完全无关。”\(^\text{[69]}\)

3月1日，《政治家报》以《原子对光的散射——新现象——加尔各答教授的发现》为标题，转载了这一新闻。\(^\text{[85]}\)拉曼于3月8日向《自然》杂志提交了一篇三段式简短报告，该文于4月21日发表。\(^\text{[86]}\)随后，他于3月22日提交了包含实际数据的论文，刊登于5月5日。\(^\text{[87]}\)3月16日，拉曼在班加罗尔召开的“南印度科学协会”会议上，以《一种新的辐射》为题，首次正式而详尽地介绍了这一发现。他的演讲发表于3月31日的《印度物理学杂志》。\(^\text{[37]}\)当天，他还将该论文的印刷版复印件一千份寄往世界各国的科学家，以便迅速传播这一重大发现。\(^\text{[38]}\)

\textbf{反响与结果}

最初，一些物理学家——尤其是法国和德国的学者——对这一发现的真实性持怀疑态度。耶拿弗里德里希·席勒大学的格奥尔格·尤斯曾致信慕尼黑大学的阿诺德·索末菲，询问道：“您认为拉曼关于液体中光学康普顿效应的研究可靠吗？……液体中散射谱线的清晰度让我感到怀疑。”索末菲随后尝试复现实验，但未能成功。\(^\text{[88]}\)1928年6月20日，柏林大学的彼得·普林斯海姆成功复现了拉曼的实验结果。他在接下来的论文中首次使用了“Ramaneffekt”（拉曼效应）和“Linien des Ramaneffekts”（拉曼谱线）等术语。\(^\text{[89][90]}\)英文版本“Raman effect”与“Raman lines”也随即被广泛采用。\(^\text{[91][23][92]}\)

除了作为一种全新的物理现象外，拉曼效应还是最早证明光具有量子性质的实验证据之一。1929年初，美国约翰斯·霍普金斯大学的罗伯特·W·伍德成为首位确认拉曼效应的美国科学家。\(^\text{[93]}\)他通过一系列实验验证后评论道：“在我看来，这一极其优美的发现，是拉曼对光散射现象长期而耐心研究的成果，也是量子理论最有力的实验证明之一。”\(^\text{[94][95]}\)以此为基础的“拉曼光谱学”由此诞生。1930年，英国皇家学会会长欧内斯特·卢瑟福在颁发休斯奖章时评价说，拉曼效应是“过去十年中实验物理领域最出色的三到四项发现之一”。\(^\text{[75]}\)

拉曼自信自己将获得诺贝尔物理学奖，但在1928年奖项授予欧文·理查森、1929年授予路易·德布罗意时，他颇感失望。他对1930年的获奖几乎笃定无疑——早在7月，他便预订了前往颁奖典礼的机票，尽管颁奖结果要到11月才公布。此后，他每天都仔细翻阅报纸，一旦没有看到诺贝尔奖的消息，便将报纸扔掉。\(^\text{[96]}\)最终，他如愿在当年获得了诺贝尔物理学奖。\(^\text{[97]}\)
\subsubsection{后期研究}
拉曼与瓦拉纳西的贝拿勒斯印度大学有着密切联系。他曾出席该校的奠基仪式\(^\text{[98]}\)，并在1916年2月5日至8日于该校举办的系列讲座中发表了关于数学与《物理学中的一些新方向》的演讲。\(^\text{[99]}\)此外，他还被聘为该校的常任访问教授。\(^\text{[100]}\)

1932年，拉曼与苏里·巴伽万塔姆共同确定了光子的自旋，进一步证实了光的量子性质。\(^\text{[101][92]}\)他与另一位学生纳根德拉·纳特合作，对声致光效应（即光被声波散射的现象）进行了理论解释，发表了一系列论文，形成了著名的“拉曼–纳特理论”。\(^\text{[102]}\)基于这一效应的调制器与开关系统，后来成为激光光学通信的重要技术基础。\(^\text{[103]}\)

他还进行了多项实验与理论研究，包括：对超声与高超声频率声波引起的光衍射的研究；\(^\text{[104][105]}\)以及X射线照射普通光下晶体所产生的红外振动效应的研究（1935年至1942年间陆续发表）。\(^\text{[106][107]}\)

1948年，拉曼通过研究晶体的光谱学行为，以一种全新的方式探讨了晶体动力学的基本问题。\(^\text{[108][109]}\)自1944年至1968年间，他持续研究了金刚石的结构与性质；\(^\text{[110][111]}\)20世纪50年代初，他又研究了多种具虹彩效应的矿物与物质的结构与光学性质，包括拉长石\(^\text{[112]}\)、珍珠长石\(^\text{[113]}\)、玛瑙\(^\text{[114]}\)、石英\(^\text{[115]}\)、蛋白石\(^\text{[116]}\)以及珍珠。\(^\text{[117]}\)此外，他还对胶体光学、电学与磁学各向异性进行了深入研究。\(^\text{[118][119]}\)进入20世纪60年代后，他的研究兴趣转向生物学领域，包括花朵的颜色机理与人类视觉的生理学特性等课题。\(^\text{[120][121][122]}\)
\subsection{个人生活}
1907年，拉曼与洛卡苏达丽·阿玛尔结婚。她是时任马德拉斯海关监督官S. 克里希纳斯瓦米·艾耶的女儿。\(^\text{[24]}\)婚礼的日期通常被记载为1907年5月6日，\(^\text{[123][124][125]}\)但据拉曼的侄孙女兼传记作者乌玛·帕拉梅斯瓦兰考证，实际日期应为1907年6月2日。\(^\text{[126][127]}\)这是一场由两人自行安排的婚姻，当时新娘年仅13岁。\(^\text{[128][43][129]}\)（不同资料对其年龄有所出入：若按其出生年份1892年来算，\(^\text{[23][124][125]}\)她应约15岁；但帕拉梅斯瓦兰确认她确为13岁，\(^\text{[130]}\)这一说法也与她的讣告相符——《当代科学》在她1980年5月22日逝世时记载其享年86岁。\(^\text{[131]}\)）她后来曾笑称，他们的婚姻与其说是因为她的音乐才华（两人初次相识时，她正在弹奏维纳琴 veena），不如说是“因为财政部门给已婚职员的额外津贴”。\(^\text{[43]}\)当时，这项“额外津贴”指的是每月为已婚官员增加150卢比的补贴。\(^\text{[24]}\)1907年夫妇二人搬到加尔各答后，曾一度被谣传改信基督教。原因是他们经常前往加尔各答的圣约翰教堂——洛卡苏达丽喜爱教堂的音乐，而拉曼则对教堂的声学效果充满兴趣。\(^\text{[42]}\)

他们育有两子：长子钱德拉塞卡尔·拉曼与次子文卡特拉曼·拉达克里希南，后者成为著名的无线电天文学家。由于思想分歧，拉曼后来与长子断绝了关系。\(^\text{[132]}\)拉曼的哥哥钱德拉塞卡拉·苏布拉马尼亚·艾耶的儿子——苏布拉马尼扬·钱德拉塞卡则在1983年获得了诺贝尔物理学奖。\(^\text{[133]}\)

拉曼一生中收集了大量具有独特光散射特性的石头、矿物与材料，这些藏品部分来自他的海外旅行，也有许多是他人赠送的礼物。\(^\text{[134]}\)他经常随身携带一台袖珍分光镜，用以研究各种样品。\(^\text{[135]}\)这些藏品连同他的分光仪，如今陈列于印度科学研究所（IISc）。\(^\text{[136][137]}\)

英国物理学家欧内斯特·卢瑟福在拉曼人生的多个关键时刻都起到了重要作用——他于1930年提名拉曼角逐诺贝尔物理学奖；同年，作为英国皇家学会会长，为他颁发了休斯奖章；1932年，又推荐他出任印度科学研究所所长一职。\(^\text{[11]}\)

拉曼对诺贝尔奖有着近乎执念的热情。他曾在加尔各答大学的一次演讲中说道：“我并不为被选为英国皇家学会院士（1924年）而感到受宠若惊。这只是一个小小的成就。若要说我真正渴望的，那就是诺贝尔奖。你们会看到，我将在五年内获得它。”\(^\text{[138]}\)他深知，如果自己真的获奖，就不能像往常那样等待诺贝尔委员会年底的公告，因为当时前往瑞典必须乘船，需要很长时间。\(^\text{[139]}\)于是，他满怀信心地在1930年7月预订了前往斯德哥尔摩的轮船票——两张，一张为自己，一张为妻子。\(^\text{[140]}\)在获得诺贝尔奖后不久的采访中，有人问他：“如果您更早发现拉曼效应，会发生什么？”他笑着回答：“那我可能就得与康普顿分享这项诺贝尔奖了，而我可不愿那样。我宁可独自获得整个奖项。”\(^\text{[141]}\)
\subsubsection{宗教观}
虽然拉曼很少谈论宗教，但他公开表示自己是一个不可知论者，\(^\text{[142]}\)并反对被称为无神论者。\(^\text{[22]}\)他的这种不可知论立场，主要受其父亲的影响——父亲崇尚赫伯特·斯宾塞、查尔斯·布拉德劳以及罗伯特·G·英格索尔等人的哲学思想。\(^\text{[143]}\)他厌恶印度教的传统仪式，\(^\text{[144]}\)但在家庭生活中并未完全摒弃这些习俗。\(^\text{[145][146]}\)此外，他也受到“不二论”哲学思想的影响。\(^\text{[147]}\)他的标志性装束是头戴传统的印度头巾（pagri，下有发髻）与佩戴印度教圣线。尽管在南印度文化中并无戴头巾的习惯，他却笑着解释道：“哦，如果我不戴头巾，我的头都会胀大。你们都这么夸我，我得用头巾把自我包住。”\(^\text{[25]}\)他甚至开玩笑地说，正是因为那顶头巾，他在首次访问英国时才获得了他人注意——尤其是J. J. 汤姆孙与卢瑟福的关注。\(^\text{[43]}\)他曾在一次公开演讲中说道：\\\\
“没有天堂、没有天界、没有地狱、没有轮回、没有再生、也没有永生。唯一真实的事实是——人出生、活着、然后死亡。因此，一个人应当好好地度过自己的一生。”\(^\text{[148]}\)

在一次与圣雄甘地和德国动物学家吉尔伯特·拉姆的友好会谈中，谈话转向宗教话题。拉曼说道：\\\\
“我来回答你（拉姆）的问题。如果真的存在上帝，我们必须在宇宙中去寻找他；如果他不在那里，那就不值得去寻找……天文学与物理学的不断新发现，似乎正在不断揭示上帝的存在。”\(^\text{[22]}\)

在生命的最后时刻，拉曼对妻子说：“我只相信人类的精神。”并嘱咐她：“我的葬礼要干净而简单——只需火化，不要那些繁琐的仪式。”\(^\text{[144]}\)
\subsection{逝世}
1970年10月底，拉曼在实验室中突发心脏骤停并倒下。他被送往医院，医生在诊断后宣告，他的生命最多只剩下四个小时。\(^\text{[149]}\)然而，拉曼顽强地又坚持了几天，并请求能回到研究所的花园中，在弟子与崇敬者们的陪伴下度过最后的时光。\(^\text{[150]}\)

在去世前两天，他对一位昔日学生说道：“不要让科学院的期刊消亡。因为它们是衡量本国科学质量的敏感指标——它们能告诉人们，科学是否真正扎根在这片土地上。”\(^\text{[43]}\)当晚，拉曼在卧室中与研究所管理委员会成员开会，讨论研究所未来的管理方向。\(^\text{[150]}\)他还嘱咐妻子，在他去世后只需为他进行简朴的火葬，不要举行任何宗教仪式。1970年11月21日清晨，拉曼因自然原因辞世，享年82岁。\(^\text{[149]}\)

拉曼逝世的消息传出后，时任印度总理英迪拉·甘地公开发表悼词：

“全国、议会以及我们每一个人都为C. V. 拉曼博士的逝世而哀悼。他是现代印度最伟大的科学家之一，也是我们国家悠久历史上最卓越的思想家之一。他的头脑如同他所研究并阐释的钻石一般晶莹剔透。他一生的事业在于揭示光的本质，而世界因他为科学所赢得的新知识而以各种方式向他致敬。”\(^\text{[151]}\)
\subsection{争议}
\subsubsection{诺贝尔奖}
\textbf{独立发现}

1928年，莫斯科国立大学的格里高利·兰兹伯格与列昂尼德·曼德尔施塔姆也独立发现了“拉曼效应”。他们的研究成果发表于当年7月的《自然科学》杂志上，\(^\text{[152]}\)并在同年8月5日至16日于萨拉托夫召开的“俄罗斯物理学家协会第六次大会”上进行了报告。\(^\text{[153]}\)1930年，二人与拉曼一同被提名诺贝尔物理学奖。然而，诺贝尔委员会最终给出了如下理由，决定仅提名拉曼：\(^\text{[154]}\)俄方并未对其发现作出独立解释，他们在论文中引用了拉曼的文章；他们只在晶体中观察到这一效应，而拉曼与克里希南则在固体、液体和气体中均观察到，从而证明了该效应的普遍性；关于拉曼光谱线与红外光谱线强度的问题已在前一年得到解释；拉曼的实验方法在分子物理的不同领域中获得了巨大成功；拉曼效应有效地用于验证分子的对称性性质，并对原子物理中与核自旋相关的问题提供了重要帮助。因此，诺贝尔委员会最终仅将拉曼一人的名字提报给瑞典皇家科学院作为获奖候选人。\(^\text{[154]}\)

然而，后来的证据表明，俄国科学家实际上更早发现了这一现象——比拉曼与克里希南早约一周。\(^\text{[155]}\)根据曼德尔施塔姆致奥列斯特·赫沃尔松的信件记载，他们于1928年2月21日便已观察到这种新的光谱线。\(^\text{[156]}\)

\textbf{克里希南的角色}

尽管克里希南是发现“拉曼效应”的主要研究者之一，但他却未被提名诺贝尔奖。\(^\text{[88]}\)最初发现这一新型散射现象的，正是他本人。\(^\text{[64]}\)在1928年有关该发现的所有科学论文中，除两篇外，其他全部由他与拉曼共同署名。此后所有的后续研究论文也均由克里希南独立撰写。\(^\text{[157][158][159]}\)然而，克里希南本人从未宣称自己理应获得诺贝尔奖。\(^\text{[160]}\)拉曼后来也承认，克里希南确实是“共同发现者”。\(^\text{[88]}\)但他与克里希南之间的关系长期紧张，甚至公开对立。克里希南晚年曾哀叹，这种敌意是“我一生中最大的悲剧”。\(^\text{[160]}\)在克里希南去世后，拉曼曾对《印度时报》的一位记者说道：“克里希南是我所见过的最大骗子，他一生都披着别人的发现而活。”\(^\text{[161:251]}\)
\subsubsection{拉曼－玻恩之争}
1933年10月至1934年3月间，马克斯·玻恩受拉曼邀请，受聘为印度科学研究所理论物理学讲师。【162】当时，玻恩因逃离纳粹德国而成为流亡学者，暂居英国剑桥圣约翰学院。\(^\text{[163]}\)自20世纪初起，玻恩便基于热学性质提出了一种晶格动力学理论。\(^\text{[164]}\)他在印度科学研究所的一次讲座中介绍了这套理论。然而，当时拉曼已发展出另一种理论，并声称玻恩的理论与实验数据相矛盾。\(^\text{[162]}\)两人之间的争论从此展开，并持续了数十年之久。\(^\text{[165][166]}\)

在这场争论中，大多数物理学家支持玻恩，\(^\text{[167]}\)因为事实证明他的理论在解释晶格行为方面更为准确。\(^\text{[162]}\)拉曼的理论虽然具有一定启发意义，但普遍被认为只具有部分适用性。\(^\text{[168]}\)然而，这场分歧不仅限于学术层面，还蔓延至个人与社会关系。玻恩后来回忆道，拉曼“或许把我当作敌人”。\(^\text{[162]}\)即便后来越来越多的证据支持玻恩的理论，拉曼仍拒绝承认其正确性。作为《当代科学》的主编，他甚至拒绝刊登任何支持玻恩理论的论文。\(^\text{[169]}\)玻恩后来多次因其在晶格理论方面的贡献被提名诺贝尔物理学奖，最终于1954年因其在量子力学统计解释方面的工作获奖。人们常将这一奖项称为他“迟来的诺贝尔奖”。\(^\text{[170]}\)
\subsubsection{与印度当局的关系}
拉曼对当时的印度总理贾瓦哈拉尔·尼赫鲁及其科学政策极为反感。他曾有一次将尼赫鲁的半身像摔在地上打碎；另一次，他更是亲手用锤子砸碎了自己的“印度国宝勋章”奖章——因为该奖项是由尼赫鲁政府授予的。【171】【172】1948年，尼赫鲁访问拉曼研究所时，拉曼当众讽刺了他。当时，工作人员在紫外光下展示了一块金属样品——一块金和一块铜。尼赫鲁误以为那块在紫外光下发出更明亮光芒的铜是金子。拉曼立刻调侃道：“总理先生，并非所有闪闪发光的东西都是金子。”\(^\text{[173]}\)

同一次访问中，尼赫鲁提出愿意为研究所提供财政资助，但拉曼断然拒绝，说道：“我绝不希望这里变成另一个政府实验室。”\(^\text{[149]}\)拉曼尤其反对政府对科研项目的集中控制，例如建立巴巴原子研究中心、国防研究与发展组织以及科学与工业研究理事会。\(^\text{[169][174]}\)他对这些机构及其领导人物——包括霍米·J·巴巴、S. S. 巴特纳加尔），甚至他昔日最得意的学生克里希南——都持敌对态度。他甚至将这类体制化科研讽刺为“尼赫鲁－巴特纳加尔效应”。\(^\text{[175][176][177]}\)1959年，拉曼提议在马德拉斯建立另一所研究机构。马德拉斯政府建议他向中央政府申请经费，但拉曼早已预见结果，他在回复时对时任马德拉斯财政与教育部长C. 苏布拉马尼安说道：“如果我把这份提案交给尼赫鲁政府，结果只会是被拒绝。”此事因此作罢。\(^\text{[174]}\)

拉曼还曾讽刺印度国大党的官员们，说他们“只是一场大Tamasha（戏剧或闹剧）”，整天讨论问题却从不付诸行动。至于印度粮食短缺问题，他向政府直言：“我们必须停止像猪一样繁殖，那问题自然就会解决。”\(^\text{[138]}\)

\textbf{印度科学院}

印度科学院的诞生，源于当时在建立类似英国皇家学会的国家级科学机构提案过程中所产生的分歧。\(^\text{[178]}\)1933年，当时印度规模最大的科学组织——印度科学大会协会计划成立一个国家科学机构，以便在科学事务上向政府提供咨询。\(^\text{[179]}\)彼时，《自然》杂志主编理查德·格雷戈里爵士访问印度时，建议作为《当代科学》主编的拉曼建立一所“印度科学院”。然而，拉曼认为该学会应当完全由印度人组成，而不是像多数人建议的那样接纳英国籍成员。他坚决表态说：“印度科学怎能在一个理事会中繁荣？其中30名成员有15人是英国人，而这15人中只有两三人真正有资格成为院士。”1933年4月1日，拉曼召集了南印度科学家召开独立会议，他与苏巴·拉奥随后正式退出印度科学大会协会。\(^\text{[180]}\)

4月24日，拉曼将新组织正式注册为“印度科学院”，提交至社团注册局。\(^\text{[179]}\)该名称原为临时称谓，计划在获得英国皇家特许后更名为“印度皇家学会”。然而，印度政府并未承认其为国家级科学机构。于是，印度科学大会协会于1935年1月7日另行创建了一个新组织——“印度国家科学研究院”，后于1970年更名为“印度国家科学院”）。\(^\text{[180]}\)该新机构由拉曼的几位主要学术对手领导，包括梅格纳德·萨哈、霍米·J·巴巴、S. S. 巴特纳加尔以及K. S. 克里希南。\(^\text{[178]}\)

\textbf{印度科学研究所}

拉曼与印度科学研究所管理层之间的关系一度恶化。他被指责在推动物理学发展时存在偏袒，而忽视了其他学科的建设。\(^\text{[162]}\)他在与同僚的相处上缺乏外交手腕。其侄子、后来的IISc所长S. 拉马塞尚曾回忆说：“拉曼走进研究所时，就像一头闯进瓷器店的公牛。”\(^\text{[144]}\)拉曼希望将物理研究提升到与西方研究机构同等的水准，但这种努力往往是以牺牲其他科学领域为代价的。\(^\text{[162]}\)马克斯·玻恩对此也曾评论说：“拉曼发现了一个懒散的地方，那里聚集着一群高薪却几乎不干活的人。”\(^\text{[144]}\)在一次理事会会议上，电气技术系教授肯尼斯·阿斯顿严厉批评了拉曼及其聘请玻恩的决定。拉曼原本打算让玻恩出任正教授。\(^\text{[25]}\)阿斯顿甚至进行了人身攻击，称玻恩是“一个被自己国家拒绝的叛徒，因此只是个二流科学家，不配成为教员，更不该领导物理系。”\(^\text{[181]}\)

1936年1月，IISc理事会成立了一个审查委员会，以评估拉曼的行为。该委员会由圣安德鲁斯大学校长兼副校长詹姆斯·欧文担任主席。委员会于3月提交报告，指控拉曼滥用资金，并“完全将研究所的重心转移到了物理学研究上”。报告还指出，理事会早在1935年11月批准的“聘任玻恩为数学物理学教授”的提案在财政上不可行。\(^\text{[162]}\)最终，理事会向拉曼提出了两个选择：从1936年4月1日起辞去研究所职务；辞去所长职务，但继续担任物理学教授。如果他拒绝选择其中之一，将被直接解职。拉曼倾向于选择第二个选项。\(^\text{[182]}\)
\subsubsection{英国皇家学会}
拉曼似乎从未对“英国皇家学会院士”这一头衔抱有太高评价。\(^\text{[138]}\)1968年3月9日，他正式提交辞呈，请求辞去院士身份。皇家学会理事会于4月4日批准了这一决定，但具体原因并未在官方记录中说明。\(^\text{[183]}\)其中一个可能的原因，是拉曼反对皇家学会仍在院士类别中使用“英国臣民”这一称谓。尤其是在印度独立之后，皇家学会内部也因这一称谓问题而产生了争议。\(^\text{[184]}\)

据苏布拉马尼扬·钱德拉塞卡回忆，《伦敦时报》曾刊登过一份皇家学会院士名单，却遗漏了拉曼的名字。拉曼因此致信时任学会会长帕特里克·布莱基特，要求解释。然而布莱基特回复称，学会对报纸的行为不负任何责任，这令拉曼深感失望。\(^\text{[161:253]}\)根据克里希南的说法，另一原因可能是拉曼曾向《英国皇家学会学报》提交论文，却收到一份否定性的评审意见。这些积怨或许共同促成了他的决定。在辞呈中，拉曼写道：“我是在仔细权衡所有相关情况后作出这一决定的。我请求贵会接受我的辞职，并将我的名字从院士名单中移除。”\(^\text{[183]}\)
\subsection{荣誉与奖项}
\begin{figure}[ht]
\centering
\includegraphics[width=6cm]{./figures/03bf4ed4b11bc654.png}
\caption{比尔拉工业与技术博物馆花园中的钱德拉塞卡拉·文卡塔·拉曼半身像} \label{fig_QDLlm_5}
\end{figure}
拉曼一生获得了众多荣誉博士学位及多个科学学会的会员称号。在印度国内，他不仅是印度科学院的创始人兼首任院长，\(^\text{[185]}\)还是孟加拉亚洲学会的院士，\(^\text{[186]}\)并自1943年起成为印度科学促进会的创始院士。\(^\text{[187]}\)1935年，他被任命为印度国家科学研究院（National Institute of Sciences of India，现为印度国家科学院 Indian National Science Academy）的创始院士。\(^\text{[186][n 1][n 2]}\)此外，他还是多个国际学术团体的成员，包括：慕尼黑德意志科学院、苏黎世瑞士物理学会、格拉斯哥皇家哲学学会、爱尔兰皇家科学院、匈牙利科学院、苏联科学院、美国光学学会名誉会员、美国矿物学会、罗马尼亚科学院、美国弦乐声学学会，以及捷克斯洛伐克科学院。\(^\text{[190]}\)

1924年，他当选为英国皇家学会院士。\(^\text{[4]}\)然而，他于1968年主动辞去这一职务，成为唯一一位辞去FRS头衔的印度科学家。\(^\text{[191]}\)1929年，他担任第16届印度科学大会的大会主席。1933年至去世前，他一直担任印度科学院的院长与主要推动者。\(^\text{[96]}\)1961年，他被选为梵蒂冈科学院的院士。\(^\text{[190]}\)
\subsubsection{奖项}
\begin{itemize}
\item 1912年：拉曼仍在印度财政部任职期间，获得了柯松研究奖。\(^\text{[192]}\)
\item 1913年：再次在财政部任职时，荣获伍德本研究奖章。\(^\text{[192]}\)
\item 1928年：获意大利罗马国家科学院颁发的马泰乌奇奖章。\(^\text{[190]}\)
\item 1930年：被授予爵士头衔。他原应入选1929年国王生日授勋名单，但因审批延误，最终由印度总督欧文勋爵在新德里的总督府（今印度总统府 Rashtrapati Bhavan）为他举行特别授勋仪式。\(^\text{[193][194]}\)
\item 1930年：荣获诺贝尔物理学奖，表彰其“在光的散射方面的研究及发现以他名字命名的效应”。\(^\text{[97]}\)他是首位获得诺贝尔科学奖的亚洲人，也是第一位获此殊荣的非白人。在他之前，印度诗人拉宾德拉纳特·泰戈尔于1913年获得诺贝尔文学奖。
\item 1930年：获英国皇家学会颁发的休斯奖章。\(^\text{[75]}\)
\item 1941年：获美国费城富兰克林学会颁发的富兰克林奖章。\(^\text{[190]}\)
\item 1954年：获印度最高平民荣誉“印度国宝勋章”，与政治家、前印度总督C. 拉贾戈帕拉查里及哲学家萨尔维帕利·拉达克里希南爵士同年获奖。\(^\text{[195][196]}\)
\item 1957年：获苏联颁发的列宁和平奖。\(^\text{[190]}\)
\end{itemize}
\subsubsection{身后荣誉与当代纪念}
\begin{figure}[ht]
\centering
\includegraphics[width=6cm]{./figures/7d3e5101060604ab.png}
\caption{} \label{fig_QDLlm_6}
\end{figure}
\begin{itemize}
\item 印度每年于2月28日庆祝“全国科学日”，以纪念拉曼于1928年发现“拉曼效应”。\(^\text{[197]}\)
\item 1971年与2009年，印度邮政分别发行了印有拉曼肖像的纪念邮票。\(^\text{[198]}\)
\item 在印度首都新德里，有一条道路以他的名字命名——C. V. Raman Marg。【\(^\text{[199]}\)
在班加罗尔东部，有一片区域被称为C. V. Raman Nagar。\(^\text{[200]}\)
\item 班加罗尔国家研讨会中心北侧的道路亦被命名为C. V. Raman Road。\(^\text{[201]}\)
\item 在金奈迈拉普尔区，也有一条C. V. Raman Road。
\item 印度科学研究所的一栋建筑以他命名，称为Raman Building。\(^\text{[202]}\)
\item 班加罗尔东部80英尺大道（80 Ft. Rd.）上的一家医院被命名为Sir C. V. Raman Hospital。\(^\text{[203]}\)
\item 在他出生地蒂鲁吉拉帕利，也有一个地区被命名为C. V. Raman Nagar。
\item 月球上一座环形山以他命名，称为Raman陨石坑。
\item 以他名字命名的C. V. Raman Global University于1997年成立。
\item 1998年，美国化学学会与印度科学促进会在加尔各答贾达普尔的印度科学促进会内，联合将拉曼的发现认定为“国际化学历史地标”。纪念碑上的铭文写道：\\\\
“在这所研究所中，C. V. 拉曼爵士于1928年发现：当一束彩色光线进入液体时，其中一部分被散射的光会呈现出不同的颜色。拉曼证明，这种被散射光的性质取决于样品的类型。科学家们很快认识到这一现象在分析与研究中的重要意义，并将其命名为‘拉曼效应’。随着现代计算机与激光技术的发展，这一方法的价值更为凸显。如今，它被广泛应用于从无损矿物鉴定到重大疾病早期检测等多个领域。因这一发现，拉曼于1930年荣获诺贝尔物理学奖。”\(^\text{[35]}\)
\item C. V. 拉曼纪念机构与大众文化中的形象位于恰蒂斯加尔邦的C. V. 拉曼大学成立于2006年。
\item 2013年11月7日，谷歌推出了纪念拉曼诞辰125周年的Google Doodle主页动画，以纪念这位伟大的物理学家。\(^\text{[204][205][206]}\)
\item 位于那格浦尔的拉曼科学中心也以他命名。\(^\text{[207]}\)
\item 此后，C. V. 拉曼大学（比哈尔分校）与C. V. 拉曼大学（坎德瓦分校）分别于2018年成立。
\end{itemize}
\subsection{大众文化中的拉曼形象}
\begin{itemize}
\item 《C. V. Raman: The Scientist and His Legacy》（〈C. V. 拉曼：科学家与他的遗产〉）由南丹·库迪亚迪执导的传记电影，于1989年上映，荣获印度国家电影奖“最佳传记片奖”。\(^\text{[208]}\)
\item 《Beyond Rainbows: The Quest & Achievement of Dr. C.V. Raman》（〈越过彩虹：拉曼的追寻与成就〉）由阿南娅·班纳吉执导的纪录片，于2004年在印度国家电视台多尔达山播出。\(^\text{[209]}\)
\item 《Rocket Boys》（〈火箭之子〉）一部在SonyLIV平台播出的印地语传记剧集。剧中拉曼一角由T.M. Karthik饰演。\(^\text{[210]}\)
\end{itemize}
\subsection{相关主题（参见）}
\begin{itemize}
\item 相干反斯托克斯拉曼光谱学
\item 《拉曼光谱学期刊》
\item 拉曼放大
\item 拉曼激光
\item 拉曼显微镜
\item 拉曼光学活性
\item 共振拉曼光谱
\item 旋转偏振相干反斯托克斯拉曼光谱
\item SHERLOC一种用于火星探测的紫外拉曼光谱仪）
\item 空间偏移拉曼光谱
\item 受激拉曼绝热过程
\item 表面增强拉曼光谱
\item 探针增强拉曼光谱
\item 透射拉曼光谱
\item X射线拉曼散射
\item 钱德拉塞卡家族
\end{itemize}
\subsection{附注}
\begin{enumerate}
\item 1970年前，印度国家科学院原名为“印度国家科学研究院”，其院士称号后缀为“FNI”。1970年改为现名后，后缀改为“FNA”。
\item 虽然拉曼于1935年当选院士，但他未完成成为正式院士所需的手续，因此其资格被视为失效。\(^\text{[188]}\)不过，学院至今仍在其院士名录中将他列为“已故院士”。\(^\text{[189]}\)
\end{enumerate}
\subsection{参考文献}
\begin{enumerate}
\item “Chandrasekhara Venkata Raman - The Mathematics Genealogy Project”。mathgenealogy.org。检索于2025年6月2日。
\item  “RAMAN Definition & Meaning”。dictionary.com。检索于2025年6月2日。
\item “C.V. Raman and the Raman Effect | Biography of Sir C.V. Raman”。acs.org，美国化学学会。检索于2024年11月15日。
\item Bhagavantam, Suri（1971）。〈Chandrasekhara Venkata Raman, 1888–1970〉，《皇家学会院士传记回忆录》，第17卷，第564–592页。doi:10.1098/rsbm.1971.0022。
\item Singh, Rajinder; Riess, Falk（1998）。〈Sir C. V. Raman and the story of the Nobel prize〉，《当代科学》，第75卷（第9期）：965–971。JSTOR 24101681。
\item Deka, Mridusmita（2020年11月7日）。〈CV Raman Birth Anniversary 2020: Interesting Facts About The Nobel Laureate〉。NDTV.com。检索于2024年7月13日。
\item Goswami, Madhusree（2021年11月7日）。〈*Lesser Known Facts About Sir CV Raman*〉。《The Logical Indian》。检索于2024年7月13日。
\item Prasar, Vigyan。〈Chandrasekhara Venkata Raman: A Legend of Modern Indian Science〉。印度政府出版。存档于2013年11月10日。检索于2013年11月7日。
\item “CV RAMAN: A Creative Mind Par Excellence”。《印度斯坦时报》，2019年7月8日。存档于2020年3月3日。检索于2020年3月8日。
\item Jayaraman, Aiyasami（1989）。《Chandrasekhara Venkata Raman: A Memoir》。班加罗尔：印度科学院。第4页。ISBN 978-81-85336-24-4。OCLC 21675106。
\item 克拉克（Clark）, Robin J. H.（2013）。〈Rayleigh, Ramsay, Rutherford and Raman – their connections with, and contributions to, the discovery of the Raman effect〉。《分析家》，第138卷（第3期）：729–734。Bibcode:2013Ana...138..729C。doi:10.1039/C2AN90124B。PMID 23236600。
\item “CV RAMAN: A Creative Mind Par Excellence”。《印度斯坦时报》。2019年7月8日。存档于2020年3月3日。检索于2020年3月8日。
\item “Remembering CV Raman on his death anniversary”。《Udayavani – ಉದಯವಾಣಿ》。存档于2021年8月12日。检索于2020年3月8日。
\item 穆克吉与穆克赫帕德亚（2018）。〈Sir Chandrasekhara Venkata Raman (1888–1970)〉，载于《加尔各答物理科学学派的历史》。新加坡：新加坡施普林格出版社，第21–76页。doi:10.1007/978-981-13-0295-4_2。ISBN 978-981-13-0294-7。
\item “本月物理史：1928年2月——拉曼散射的发现”。美国物理学会新闻档案，2009年第18卷第2期。存档于2013年5月23日。
\item 拉曼（Raman, C. V.）（1906）。〈LV. Unsymmetrical diffraction-bands due to a rectangular aperture〉。《伦敦、爱丁堡及都柏林哲学杂志与科学期刊》，第12卷（第71期）：494–498。doi:10.1080/14786440609463564。存档于2020年10月31日。检索于2020年3月10日。
\item 《1930年诺贝尔物理学奖：文卡塔·拉曼爵士》。诺贝尔奖官方网站。存档于2019年2月27日。
\item “About C. V. Raman: Life, Achievements and Paper Publications”。印度中央邦印多尔。2020年2月13日。存档于2020年2月15日。检索于2020年2月20日。
\item 拉曼（Raman, C. V.）（1907）。〈LVIII. The curvature method of determining the surface-tension of liquids〉。《伦敦、爱丁堡及都柏林哲学杂志与科学期刊》，第14卷（第83期）：591–596。doi:10.1080/14786440709463720。存档于2020年10月31日。检索于2020年3月10日。
\item 瑞利勋爵（1907）。〈LLX. On the relation of the sensitiveness of the ear to pitch, investigated by a new method〉。《伦敦、爱丁堡及都柏林哲学杂志与科学期刊》，第14卷（第83期）：596–604。doi:10.1080/14786440709463721。存档于2020年10月31日。检索于2020年8月24日。
\item 辛格与里斯（2004）。〈The Nobel Laureate Sir Chandrasekhara Venkata Raman FRS and His Contacts with the British Scientific Community in a Social and Political Context〉。《伦敦皇家学会笔记与记录》，第58卷（第1期）：47–64。doi:10.1098/rsnr.2003.0224。JSTOR 4142032。S2CID 144713213。
\item 贾亚拉曼（1989）。同上，第5页。OCLC 21675106。
\item 辛格（2002）。〈C.V. Raman and the Discovery of the Raman Effect〉。《物理学视野》，第4卷（第4期）：399–420。Bibcode:2002PhP.....4..399S。doi:10.1007/s000160200002。S2CID 121785335。
\item 贾亚拉曼（1989）。同上，第8页。OCLC 21675106。
\item 班纳吉（2014）。〈C. V. Raman and Colonial Physics: Acoustics and the Quantum〉。《物理学视野》，第16卷（第2期）：146–178。Bibcode:2014PhP....16..146B。doi:10.1007/s00016-014-0134-8。S2CID 121952683。
\item 比瓦斯（2010）。〈Indian Association for the Cultivation of Science: A Nation's Dream, 1869–1947〉。载于达斯古普塔（编），《科学与现代印度：一个制度史，约1784–1947》。德里：培生出版社（Pearson），第69–116页。ISBN 978-93-325-0294-9。OCLC 895913622。
\item 拉曼（Raman, C. V.）（1907）。〈Newton's rings in polarised light〉。《自然》，第76卷（第1982期）：637。Bibcode:1907Natur..76..637R。doi:10.1038/076637b0。S2CID 4035854。存档于2020年10月31日。检索于2020年3月30日。
\item 巴苏（2016）。《C.V. Raman的一生与时代》。新德里：Prabhat Prakashan出版社，第22–23页。ISBN 978-81-8430-362-9。存档于2021年8月12日。检索于2020年6月3日。
\item 《C.V. Raman: A Pictorial Biography》。印度科学院出版。班加罗尔，1988年。第3页。ISBN 978-81-85324-07-4。检索于2018年2月26日。
\item 穆克吉与穆克赫帕德亚（2018）。《加尔各答物理科学学派的历史》。新加坡：施普林格出版社，第30页。ISBN 978-981-13-0295-4。OCLC 1042158966。
\item 穆克吉。同上，第31页。OCLC 1042158966。
\item “C.V. Raman”。美国光学学会（OSA），华盛顿特区。2013年6月12日。存档于2021年8月12日。检索于2020年3月8日。
\item 布兰皮德（1986）。〈Pioneer Scientists in Pre-Independence India〉。《今日物理》，第39卷（第5期）：36–44。Bibcode:1986PhT....39e..36B。doi:10.1063/1.881025。
\item 贾亚拉曼与兰达斯（1988）。〈Chandrasekhara Venkata Raman〉。《今日物理》，第41卷（第8期）：56–64。Bibcode:1988PhT....41h..56J。doi:10.1063/1.881128。
\item “C.V. Raman: The Raman Effect – Landmark”。美国化学学会。存档于2020年3月4日。检索于2020年3月8日。
\item “Indian Journal of Physics”。1926年创刊。存档于2018年3月8日。检索于2018年3月12日。
\item 拉曼（1928）。〈A new radiation〉。《印度物理学杂志》，第2卷：387–398。存档于2019年5月20日。检索于2018年3月12日。
\item 贾亚拉曼（1989）。*同上*，第30页。OCLC 21675106。
\item “Indian Association for the Cultivation of Science (1876–)”。印度科学促进会。存档于2017年7月17日。检索于2018年3月5日。
\item 雷迪与古塔尔。“Indian Institute of Science”。《The Conversation》。存档于2020年10月20日。检索于2020年3月18日。
\item 帕拉梅斯瓦兰（2011）。《C.V. Raman: A Biography》（《C.V. 拉曼传记》）。新德里：企鹅印度出版社。第190页。ISBN 978-0-14-306689-7。OCLC 772714846。存档于2021年7月27日。
\item “About us”。TCM有限公司官方网站。存档于2014年3月1日。检索于2013年11月6日。
\item 马斯卡雷尼亚斯（2013）。〈Sir C.V. Raman – Icon of Indian Science〉。《科学记者》，第50卷（第11期）：21–25。存档于2020年10月29日。检索于2020年3月15日。
\item 拉曼（1916）。〈XLIII. On the "wolf-note" in bowed stringed instruments〉。《伦敦、爱丁堡及都柏林哲学杂志与科学期刊》，第32卷（第190期）：391–395。doi:10.1080/14786441608635584。
\item 拉曼（1916）。〈On the "Wolf-note" of the Violin and 'Cello〉。《自然》，第97卷（第2435期）：362–363。Bibcode:1916Natur..97..362R。doi:10.1038/097362a0。S2CID 3966106。存档于2020年10月31日。检索于2020年4月30日。
\item 拉曼（1918）。〈On the mechanical theory of the vibrations of bowed strings and of musical instruments of the violin family, with experimental verification of the results – Part I〉（PDF）。《印度科学促进会通报》，第15卷：1–158。PDF存档于2020年10月23日。检索于2020年3月10日。
\item 拉曼（Raman, C. V.）（1921）。〈On some Indian stringed instruments〉（PDF）。《印度科学促进会会刊》，第7卷：29–33。PDF存档于2014年1月2日。检索于2020年3月10日。
\item 拉曼与阿舒托什·德伊（Dey, Ashutosh）（1920）。〈X. On the sounds of splashes〉。《伦敦、爱丁堡及都柏林哲学杂志与科学期刊》，第39卷（第229期）：145–147。doi:10.1080/14786440108636021。
\item 拉曼（1920）。〈Experiments with mechanically-played violins〉（PDF）。《印度科学促进会会刊》，第6卷：19–36。PDF存档于2020年10月23日。检索于2020年3月10日。
\item 拉曼与库马尔（1920）。〈Musical Drums with Harmonic Overtones〉。《自然》，第104卷（第2620期）：500。Bibcode:1920Natur.104..500R。doi:10.1038/104500a0。S2CID 4159476。存档于2020年10月31日。检索于2020年8月24日。

\item 拉曼与西瓦卡利·库马尔（1920）。〈Musical drums with harmonic overtones*〉。《自然》，第104卷（第2620期）：500。Bibcode:1920Natur.104..500R。doi:10.1038/104500a0。S2CID 4159476。存档于2020年10月31日。检索于2020年8月24日。
\item 拉曼与班纳吉（1920）。〈On Kaufmann's theory of the impact of the pianoforte hammer〉。《伦敦皇家学会学报A辑：数学与物理论文集》，第97卷（第682期）：99–110。Bibcode:1920RSPSA..97...99R。doi:10.1098/rspa.1920.0016。
\item 拉曼与萨瑟兰（1921）。〈Whispering-Gallery Phenomena at St. Paul's Cathedral〉。《自然》，第108卷（第2706期）：42。Bibcode:1921Natur.108...42R。doi:10.1038/108042a0。S2CID 4126913。
\item 拉曼（1922）。〈On whispering galleries〉（PDF）。《印度科学促进会通报》，第7卷：159–172。PDF存档于2013年12月12日。检索于2015年1月7日。
\item 班纳吉（（2014）。〈C. V. Raman and Colonial Physics: Acoustics and the Quantum〉。《物理学视野》，第16卷（第2期）：146–178。Bibcode:2014PhP....16..146B。doi:10.1007/s00016-014-0134-8。S2CID 121952683。
\item 拉曼（1919）。〈LVI. The scattering of light in the refractive media of the eye〉。《伦敦、爱丁堡及都柏林哲学杂志与科学期刊》，第38卷（第227期）：568–572。doi:10.1080/14786441108635985。存档于2020年10月31日。检索于2020年3月10日。
\item 佚名）（2009）。〈This Month in Physics History: February 1928: Raman scattering discovered〉。《美国物理学会新闻》，第12卷（第2期）：在线版。存档于2020年3月6日。检索于2020年3月10日。
\item 布坎南（1910）。〈Colour of the Sea〉。《自然》，第84卷（第2125期）：87–89。Bibcode:1910Natur..84...87B。doi:10.1038/084087a0。
\item 巴恩斯（1910）。〈Colour of Water and Ice〉。《自然》，第83卷（第2111期）：188。Bibcode:1910Natur..83..188B。doi:10.1038/083188a0。S2CID 3943242。存档于2020年10月31日。检索于2020年3月16日。
\item 瑞利（（1910）。〈Colours of Sea and Sky〉。《自然》（，第83卷（第2106期）：48–50。Bibcode:1910Natur..83...48..。doi:10.1038/083048a0。
\item 瑞利勋爵（1899）。〈XXXIV. On the transmission of light through an atmosphere containing small particles in suspension, and on the origin of the blue of the sky〉（《含悬浮微粒大气中光的传播及天空蓝色的起源》）。《伦敦、爱丁堡及都柏林哲学杂志与科学期刊》，第47卷（第287期）：375–384。doi:10.1080/14786449908621276。存档于2020年10月31日。检索于2020年3月16日。
\item 施特特费尔德、麦肯纳与帕特尔（2016）。〈Dynamic light scattering: a practical guide and applications in biomedical sciences〉（《动态光散射：实用指南及其在生物医学中的应用》）。《生物物理评论》，第8卷（第4期）：409–427。doi:10.1007/s12551-016-0218-6。PMC 5425802。PMID 28510011。
\item 拉曼（1921）。〈The Colour of the Sea〉（《海水的颜色》）。《自然》，第108卷（第2716期）：367。Bibcode:1921Natur.108..367R。doi:10.1038/108367a0。S2CID 4064467。
\item 马利克（2000）。〈The Raman Effect and Krishnan's Diary〉（《拉曼效应与克里希南的日记》）。《伦敦皇家学会记录与文献》，第54卷（第1期）：67–83。doi:10.1098/rsnr.2000.0097。JSTOR 532059。S2CID 143485844。
\item 拉曼（1922）。〈On the molecular scattering of light in water and the colour of the sea〉（《水中光的分子散射与海水的颜色》）。《伦敦皇家学会学报A辑：数学与物理论文集》，第101卷（第708期）：64–80。Bibcode:1922RSPSA.101...64R。doi:10.1098/rspa.1922.0025。
\item 拉曼纳森（1923）。〈LVIII. On the colour of the sea〉（《海水的颜色》）。《伦敦、爱丁堡及都柏林哲学杂志与科学期刊》，第46卷（第273期）：543–553。doi:10.1080/14786442308634277。
\item 拉曼纳森（1925年3月1日）。〈The Transparency and Color of the Sea〉（《海水的透明度与颜色》）。《物理评论》，第25卷（第3期）：386–390。Bibcode:1925PhRv...25..386R。doi:10.1103/PhysRev.25.386。
\item 布劳恩与斯米尔诺夫（1993）。〈Why is water blue?〉（《为什么水是蓝色的？》）（PDF）。《化学教育杂志》），第70卷（第8期）：612–614。Bibcode:1993JChEd..70..612B。doi:10.1021/ed070p612。PDF存档于2019年12月1日。检索于2021年7月27日。
\item 辛格（2018年3月1日）。〈The 90th Anniversary of the Raman Effect〉（《拉曼效应九十周年纪念》）（PDF）。《印度科学史杂志》，第53卷（第1期）：50–58。doi:10.16943/ijhs/2018/v53i1/49363。PDF存档于2020年8月25日。检索于2020年3月17日。
\item 克里希南（1925）。〈LXXV. On the molecular scattering of light in liquids〉（《液体中光的分子散射》）。《伦敦、爱丁堡及都柏林哲学杂志与科学期刊》，第50卷（第298期）：697–715。doi:10.1080/14786442508634789。
\item 马利克（2000）。〈The Raman effect and Krishnan's diary〉（《拉曼效应与克里希南的日记》）。《伦敦皇家学会记录与文献》，第54卷（第1期）：67–83。doi:10.1098/rsnr.2000.0097。S2CID 143485844。存档于2019年10月24日。检索于2020年12月13日。
\item 拉曼与克里希南（1927）。〈Magnetic double-refraction in liquids. part I.—benzene and its derivatives〉（《液体中的磁致双折射现象——第一部分：苯及其衍生物》）。《伦敦皇家学会学报A辑：数学与物理论文集》，第113卷（第765期）：511–519。Bibcode:1927RSPSA.113..511R。doi:10.1098/rspa.1927.0004。
\item 康普顿（1923年5月）。〈A Quantum Theory of the Scattering of X-Rays by Light Elements〉（《X射线被轻元素散射的量子理论》）。《物理评论》（，第21卷（第5期）：483–502。Bibcode:1923PhRv...21..483C。doi:10.1103/PhysRev.21.483。
\item 拉曼（1927）。〈Thermodynamics, Wave-theory, and the Compton Effect〉（《热力学、波动理论与康普顿效应》）。《自然》，第120卷（第3035期）：950–951。Bibcode:1927Natur.120..950R。doi:10.1038/120950a0。S2CID 29489622。
\item 拉姆达斯（1973）。〈Dr. C. V. Raman (1888–1970), Part II〉（《C. V. 拉曼博士（1888–1970），第二部分》）。《物理教育杂志》，第1卷（第2期）：2–18。存档于2020年10月31日。检索于2020年3月12日。
\item 辛格（2002）。〈C. V. Raman and the Discovery of the Raman Effect〉（《C. V. 拉曼与拉曼效应的发现》）。《物理学视野》，第4卷（第4期）：399–420。Bibcode:2002PhP.....4..399S。doi:10.1007/s000160200002。S2CID 121785335。
\item 克里希南与拉曼（1928）。〈The Negative Absorption of Radiation〉（《辐射的负吸收》）。《自然》，第122卷（第3062期）：12–13。Bibcode:1928Natur.122...12R。doi:10.1038/122012b0。ISSN 1476-4687。S2CID 4071281。
\item 拉曼与克里希南（Krishnan, K. S.）（1928）。〈A new type of secondary radiation〉（《一种新的二次辐射类型》）。《自然》，第121卷（第3048期）：501–502。Bibcode:1928Natur.121..501R。doi:10.1038/121501c0。S2CID 4128161。
\item “How is C. V. Raman connected with The National Science Day*”（《C.V. 拉曼与全国科学日的关系》）。《商业内幕》，2020年2月28日。存档于2020年4月15日。检索于2021年2月26日。
\item 朗（1988）。〈Early history of the Raman effect〉（《拉曼效应的早期历史》）。《国际物理化学评论》，第7卷（第4期）：317–349。Bibcode:1988IRPC....7..317L。doi:10.1080/01442358809353216。存档于2021年8月12日。检索于2021年7月27日。
\item 辛格（2018）。〈How costly was Raman's equipment for the discovery of Raman effect?〉（《拉曼效应的发现设备到底有多昂贵？》）（PDF）。《印度科学史杂志》，第53卷（第4期）：68–73。doi:10.16943/ijhs/2018/v53i4/49527。PDF存档于2021年7月28日。检索于2021年7月27日。
\item 《Raman Effect Visualised》（《拉曼效应可视化》）。YouTube，2008年10月8日。存档于2014年5月22日。检索于2014年5月15日。
\item 文卡塔拉曼（1988）。《Journey into Light: Life and Science of C. V. Raman》（《走进光明：C.V.拉曼的生命与科学》）。牛津大学出版社。ISBN 978-81-85324-00-5。
\item 梅赫拉与雷兴贝格（2001）。*The Historical Development of Quantum Theory, Volume 6 Part 1》（《量子理论的历史发展》第六卷第一部分）。纽约：施普林格出版社。第360页。ISBN 978-0-387-95262-8。OCLC 76255200。
\item 朗（2008）。〈80th Anniversary of the discovery of the Raman effect: a celebration〉（《拉曼效应发现八十周年纪念》）。《拉曼光谱学杂志》，第39卷（第3期）：316–321。Bibcode:2008JRSp...39..316L。doi:10.1002/jrs.1948。
\item 拉曼1928）。〈A Change of Wave-length in Light Scattering〉（《光散射中的波长变化》）。《自然》，第121卷（第3051期）：619。Bibcode:1928Natur.121..619R。doi:10.1038/121619b0。S2CID 4102940。
\item 拉曼与克里希南（1928）。〈The Optical Analogue of the Compton Effect〉（《康普顿效应的光学类比》）。《自然》，第121卷（第3053期）：711。Bibcode:1928Natur.121..711R。doi:10.1038/121711a0。S2CID 4126899。
\item 辛格（2008）。〈80 Years Ago – the Discovery of the Raman Effect at the Indian Association for the Cultivation of Science, Kolkata, India〉（《八十年前——拉曼效应在加尔各答印度科学促进会的发现》）（PDF）。《印度物理学杂志》，第82卷：987–1001。PDF原文存档于2020年6月6日。检索于2020年3月8日。
\item 普林斯海姆（1928）。〈Der Ramaneffekt, ein neuer von C. V. Raman entdeckter Strahlungseffekt〉（《拉曼效应：C.V.拉曼发现的一种新辐射现象》）（德文）。《自然科学》，第16卷（第31期）：597–606。Bibcode:1928NW.....16..597P。doi:10.1007/BF01494083。S2CID 30433182。
\item 卡雷利、普林斯海姆与罗森（Rosen, B.）（1928）。〈Über den Ramaneffekt an wässerigen Lösungen und über den Polarisationszustand der Linien des Ramaneffekts〉（《关于水溶液中的拉曼效应及其谱线偏振状态》）（德文）。《物理学杂志》，第51卷（第7–8期）：511–519。Bibcode:1928ZPhy...51..511C。doi:10.1007/BF01327842。S2CID 119516536。
\item 拉姆达斯（1928）。〈The Raman Effect and the Spectrum of the Zodiacal Light〉（《拉曼效应与黄道光谱》）。《自然》，第122卷（第3063期）：57。doi:10.1038/122057a0。S2CID 4092715。
\item 布兰德（1989年1月31日）。〈The discovery of the Raman effect〉（《拉曼效应的发现》）。《伦敦皇家学会记录与文献》，第43卷（第1期）：1–23。doi:10.1098/rsnr.1989.0001。S2CID 120964978。
\item 伍德（1929）。〈The Raman Effect with Hydrochloric Acid Gas: the 'Missing Line.'〉（《盐酸气体中的拉曼效应：“消失的谱线”》）。《自然》，第123卷（第3095期）：279。Bibcode:1929Natur.123Q.279W。doi:10.1038/123279a0。S2CID 4121854。
\item 伍德（1933）。〈Raman Spectrum of Heavy Water (By Cable)〉（《重水的拉曼光谱（电报通讯）》）。《自然》，第132卷（第3347期）：970。Bibcode:1933Natur.132..970W。doi:10.1038/132970b0。S2CID 4092727。
\item 《C.V. Raman The Raman Effect – Landmark》（《C.V.拉曼与拉曼效应——里程碑》）。美国化学学会。存档于2020年3月4日。检索于2020年3月11日。
\item 文卡塔拉曼（1995）。《Raman and His Effect》（《拉曼与他的效应》）。Orient Blackswan出版社，第50页。ISBN 978-81-7371-008-7。原文存档于2021年2月7日。检索于2016年7月28日。
\item 《The Nobel Prize in Physics 1930》（《1930年诺贝尔物理学奖》）。诺贝尔基金会。存档于2014年10月11日。检索于2008年10月9日。
\item 辛格（2013年11月8日）。〈BHU preserves CV Raman's association with university〉（《贝拿勒斯印度大学保存C.V.拉曼与该校的渊源》）。《印度时报》。存档于2013年11月28日。检索于2015年6月17日。
\item 德维迪（2011）。〈Madan Mohan Malaviya and Banaras Hindu University〉（《马丹·莫汉·马拉维亚与贝拿勒斯印度大学》）（PDF）。《当代科学》，第101卷（第8期）：1091–1095。PDF存档于2015年6月17日。检索于2015年6月17日。
\item 普拉卡什（2014年5月20日）。《Vision for Science Education》（《科学教育的愿景》）。Allied Publishers出版社，第45页。ISBN 978-81-8424-908-8。
\item 拉曼与巴加万塔姆（1932）。〈Experimental Proof of the Spin of the Photon〉（《光子自旋的实验验证》）。《自然》，第129卷（第3244期）：22–23。Bibcode:1932Natur.129...22R。doi:10.1038/129022a0。hdl:10821/664。S2CID 4064852。存档于2020年3月16日。检索于2020年3月18日。
\item 拉曼与纳坚德拉·纳特〈The diffraction of light by high-frequency sound waves. Part I〉（《高频声波对光的衍射（第一部分）》），发表于1935年《印度科学院院刊》，原文已于2014年12月13日存档于互联网档案馆。
\item 程启祥、梅萨姆·巴哈多里、马德琳·格里克、塞巴斯蒂安·鲁姆利与凯伦·伯格曼（2018）。〈Recent advances in optical technologies for data centers: a review〉（《数据中心光学技术的最新进展综述》）。《光学》，第5卷（第11期）：1354–1370。Bibcode:2018Optic...5.1354C。doi:10.1364/OPTICA.5.001354。存档于2021年8月12日。检索于2021年7月27日。
\item 拉曼与纳坚德拉·纳特（1935）。〈The diffraction of light by high frequency sound waves: Part I.〉（《高频声波对光的衍射：第一部分》）。《印度科学院院刊A辑》，第2卷（第4期）：406–412。doi:10.1007/BF03035840。S2CID 198141323。存档于2021年8月12日。检索于2021年7月27日。
\item 拉曼（1942）。〈The nature of the liquid state〉（《液态的本质》）。《当代科学》，第11卷（第8期）：303–310。JSTOR 24213500。存档于2021年7月27日。检索于2021年7月27日。
\item 拉曼（1941）。〈Crystals and photons〉（《晶体与光子》）。《印度科学院院刊A辑》，第13卷（第1期）：1–8。doi:10.1007/BF03052526。S2CID 198142013。存档于2021年8月12日。检索于2021年7月27日。
\item 拉曼（1942）。〈Reflexion of X-Rays with Change of Frequency〉（《X射线的反射与频率变化》）。《自然》，第150卷（第3804期）：366–369。Bibcode:1942Natur.150..366R。doi:10.1038/150366a0。S2CID 222373。存档于2021年7月27日。检索于2021年7月27日。
\item 拉曼（1948）。〈Dynamic X-ray reflections in crystals〉（《晶体中X射线的动态反射》）（PDF）。《当代科学》，第17卷：65–75。PDF存档于2021年7月27日。检索于2021年7月27日。
\item 拉曼（1948）。〈X-Rays and the Eigen-Vibrations of Crystal Structures〉（《X射线与晶体结构的本征振动》）。《自然》，第162卷（第4105期）：23–24。Bibcode:1948Natur.162...23R。doi:10.1038/162023b0。PMID 18939227。S2CID 4073206。存档于2021年7月27日。检索于2021年7月27日。
\item 拉曼（1943）。〈The structure and properties of diamond〉（《金刚石的结构与性质》）（PDF）。《当代科学》，第12卷（第1期）：33–42。JSTOR 24208172。PDF存档于2021年7月27日。检索于2021年7月27日。
\item 拉曼（1968）。〈The diamond: Its structure and properties〉（《金刚石：其结构与性质》）。《印度科学院院刊·A辑》，第 67 卷（第 5 期）：231–246。doi:10.1007/BF03049362。S2CID 91751475。
\item 拉曼与贾亚拉曼（A. Jayaraman）（1950）。〈The structure of labradorite and the origin of its iridescence〉（《拉长石的结构及其虹彩起源》）。《印度科学院院刊·A辑》，第 32 卷（第 1 期）：1–16。doi:10.1007/BF03172469。S2CID 128235557。
\item 拉曼与贾亚拉曼（1953）。〈The diffusion haloes of the iridescent feldspars〉（《虹彩长石的扩散晕环》）。《印度科学院院刊·A辑》，第 37 卷（第 1 期）：1–10。doi:10.1007/BF03052851。S2CID 128924627。原文存档于 2021 年 8 月 12 日。检索于 2021 年 7 月 27 日。
\item 拉曼与贾亚拉曼（1953）。〈The structure and optical behaviour of iridescent agate〉（《虹彩玛瑙的结构与光学性质》）。《印度科学院院刊·A辑》，第 38 卷（第 3 期）：199–206。doi:10.1007/BF03045221。S2CID 198139210。
\item 拉曼与贾亚拉曼（1954）。〈X-ray study of fibrous quartz〉（《纤维状石英的X射线研究》）。《印度科学院院刊·A辑》，第 40 卷（第 3 期）：107。doi:10.1007/BF03047162。S2CID 135143703。原文存档于 2021 年 8 月 12 日。检索于 2021 年 7 月 27 日。
\item 拉曼与贾亚拉曼（1953）。〈The structure and optical behaviour of iridescent opal〉（《虹彩蛋白石的结构与光学性质》）。《印度科学院院刊·A辑》，第 38 卷（第 5 期）：343–354。doi:10.1007/BF03045242。S2CID 198141813。原文存档于 2021 年 8 月 12 日。检索于 2021 年 7 月 27 日。
\item 拉曼与克里希纳穆尔蒂（1954）。〈The structure and optical behaviour of pearls〉（《珍珠的结构与光学性质》）。《印度科学院院刊·A辑》，第 39 卷（第 5 期）：215–222。doi:10.1007/BF03047140。S2CID 91310831。原文存档于 2021 年 8 月 12 日。检索于 2021 年 7 月 27 日。
\item 克里希南（（1948）。〈Sir C. V. Raman and crystal physics〉（《C. V. 拉曼与晶体物理学》）。《印度科学院院刊·A辑》，第 28 卷（第 5 期）：258。doi:10.1007/BF03170788。S2CID 172445086。原文存档于 2021 年 8 月 12 日。检索于 2021 年 7 月 27 日。
\item 拉曼与维斯瓦纳森（1955）。〈The theory of the propagation of light in polycrystalline media〉（《多晶介质中光传播理论》）。《印度科学院院刊·A辑》，第 41 卷（第 2 期）：37–44。doi:10.1007/BF03047170。S2CID 59436273。原文存档于 2021 年 8 月 12 日。检索于 2021 年 7 月 27 日。
\item 拉曼（1960）。〈The perception of light and colour and the physiology of vision〉（《光与色的感知及视觉生理学》）。《印度科学院院刊·B辑》，第 52 卷（第 6 期）：253–264。doi:10.1007/BF03047051。S2CID 170662624。
\item 拉曼（1962）。〈The role of the retina in vision〉（《视网膜在视觉中的作用》）。《印度科学院院刊·B辑》，第56卷（第2期）：77–87。doi:10.1007/BF03051587。hdl:2289/1702。JSTOR 24059723。S2CID 83159946。原文存档于2021年8月12日。检索于2021年8月12日。
\item 拉曼（1970）。〈The florachromes: Their chemical nature and spectroscopic behaviour〉（《花色素：其化学性质与光谱特性》）。《印度科学院院刊·B辑》，第72卷（第1期）：1–23。doi:10.1007/BF03049697。S2CID 97395288。原文存档于2021年8月12日。检索于2021年8月12日。
\item 《Raman, Sir (Chandrashekhara) Venkata》。《牛津国家传记辞典》，牛津大学出版社，2004年。检索于2013年10月6日。原文存档于2014年3月8日。
\item “Lokasundari Ammal”。myheritage.com。检索于2023年6月8日。
\item “C. V. Raman, Nobel Prize in Physics 1930”。Geni。1888年11月7日。检索于2023年6月8日。
\item 拉维）（2014年7月21日）。〈The Raman wife effect: lively recollections〉（《“拉曼之妻效应”：生动的回忆》）。《印度教徒报》。ISSN 0971-751X。检索于2023年6月8日。
\item 帕拉梅斯瓦兰（2011）。前引书，企鹅印度出版社，第39页。ISBN 9780143066897。
\item 米勒与考夫曼（1989）。〈C. V. Raman and the discovery of the Raman effect〉（《C.V.拉曼与拉曼效应的发现》）。《化学教育杂志》，第66卷（第10期）：795–801。Bibcode:1989JChEd..66..795M。doi:10.1021/ed066p795。ISSN 0021-9584。
\item 印度科学院（IAS）（1988）。前引书。印度科学院出版，第2页。ISBN 9788185324074。
\item 帕拉梅斯瓦兰（2011）。前引书，企鹅印度出版社，第33–34页。ISBN 9780143066897。
\item “讣告”。《当代科学》，第49卷（第11期）：425，1980年。ISSN 0011-3891。JSTOR 24083389。
\item 帕拉梅斯瓦兰与拉曼（2011）。《C. V. Raman: a biography》（《C.V.拉曼传记》）。新德里：企鹅图书出版社。ISBN 978-0-14-306689-7。
\item “S Chandrasekhar: Why Google honours him”（《S.钱德拉塞卡：为何谷歌纪念他》）。半岛电视台。2017年10月19日。原文存档于2018年9月1日。检索于2017年10月18日。
\item 《Periodic Videos》（2015年1月28日），〈Diamonds, Pearls and Atomic Bomb Stones – Periodic Table of Videos〉（《钻石、珍珠与原子弹之石——元素周期表视频》），原文存档于2021年8月12日，检索于2018年11月12日。
\item 《Periodic Videos》（2015年1月28日），〈Special Spectroscope – Periodic Table of Videos〉（《特制光谱仪——元素周期表视频》），原文存档于2019年2月11日，检索于2018年11月12日。
\item 普瓦亚（2020年2月7日）。〈A curious mind illuminated〉（《被好奇心点亮的头脑》）。《德干先锋报》。原文存档于2021年7月27日。检索于2021年7月27日。
\item 库尔卡尼（2021年6月20日）。〈How Archives in Bangalore Tell Stories about Science – Connect with IISc〉（《班加罗尔的档案如何讲述科学的故事——与印度科学研究所的联系》）。《Connect》（印度科学研究所期刊）。原文存档于2021年7月27日。检索于2021年7月27日。
\item 萨特扬（2003年7月5日）。〈The Raman Effect〉（《拉曼效应》）。《印度展望》。原文存档于2020年10月25日。检索于2020年3月15日。
\item 克拉克与克拉克（2011）。〈Rutherford and Raman – Nobel laureates who had difficult early journeys to success〉（《卢瑟福与拉曼——早年经历坎坷的诺贝尔奖得主》）。《拉曼光谱学杂志》，第42卷（第12期）：2173–2178。doi:10.1002/jrs.3061。原文存档于2021年8月12日。检索于2020年12月13日。
\item 品扎鲁与基弗（2018）。载于托波尔斯基、迪英）与霍尔里歇尔（编），〈Raman's Discovery in Historical Context〉（《拉曼的发现及其历史语境》），见《共聚焦拉曼显微镜》，收录于《表面科学丛书》，第66卷，瑞士楚格：施普林格国际出版公司，第3–21页。doi:10.1007/978-3-319-75380-5_1。ISBN 978-3-319-75378-2。检索于2020年12月13日。

\item 贾亚拉曼（Aiyasami Jayaraman）（1989）。*前引书（Op. cit.）*。企鹅印度出版社，第31页。ISBN 9780143066897。

\item 马斯卡雷尼亚斯（K. Smiles Mascarenhas）（2013）。〈*Sir C.V. Raman – Icon of Indian Science*〉（《C.V.拉曼爵士——印度科学的象征》）。《科学记者》（*Science Reporter*），第50卷（第11期）：21–25。原文存档于2020年10月29日。检索于2020年3月15日。

\item 帕拉梅斯瓦兰（Uma Parameswaran）（2011）。*前引书（Op. cit.）*。企鹅印度出版社，第5、8页。ISBN 9780143066897。

\item 拉马塞尚（S. Ramaseshan）（1988）。〈*The portrait of a scientist – C. V. Raman*〉（《科学家的肖像——C.V.拉曼》）。《当代科学》（*Current Science*），第57卷（第22期）：1207–1220。JSTOR 24091067。

\item 穆克吉（Shantanu Mukharji）（2017年11月8日）。〈*Why it's important to keep the memories of CV Raman alive*〉（《为什么保持C.V.拉曼的记忆如此重要》）。*[www.dailyo.in*。原文存档于2021年8月12日。检索于2020年3月18日。](http://www.dailyo.in*。原文存档于2021年8月12日。检索于2020年3月18日。)

\item 帕拉梅斯瓦兰（Uma Parameswaran）（2011）。*前引书（Op. cit.）*。企鹅印度出版社，第8页。ISBN 9780143066897。

\item 帕拉梅斯瓦兰（Uma Parameswaran）（2011）。*前引书（Op. cit.）*。企鹅印度出版社，第162页。ISBN 9780143066897。

\item 辛格（Rajinder Singh）（2010）。〈*Letters to the Editor: Indian scientists vs. science and religion*〉（《致编辑的信：印度科学家、科学与宗教之争》）（PDF）。《科学与文化》（*Science and Culture*），第76卷（第7–8期）：206。PDF存档于2020年7月12日。检索于2020年3月18日。

\item 库尔卡尼（Pavan Kulkarni）（2015年11月20日）。〈*The last years: Raman's meeting with Nehru and more*〉（《最后的岁月：拉曼与尼赫鲁的会面及其后》）。《公民事务》（*Citizen Matters*）。原文存档于2020年2月19日。检索于2020年3月15日。

\item 印度科学院（IAS）（1988）。*前引书（Op. cit.）*。印度科学院出版，第177页。ISBN 9788185324074。

\item 甘地，英迪拉（Indira Gandhi）（1975）。《*Selected Speeches of Indira Gandhi: The Years of Endeavour, August 1969 – August 1972*》（《英迪拉·甘地演讲选集：奋斗的岁月，1969年8月—1972年8月》）。新德里：印度政府新闻与广播部出版司，第804页。ISBN 978-0-940500-97-6。OCLC 2119197。

\item 兰茨伯格（G.）与曼德尔施塔姆（L.）（1928）。〈*Eine neue Erscheinung bei der Lichtzerstreuung in Krystallen*〉（《晶体中光散射的一种新现象》）。《自然科学》（*Die Naturwissenschaften*，德文），第16卷（第28期）：557–558。Bibcode:1928NW.....16..557..。doi:10.1007/BF01506807。S2CID 22492141。

\item 乌萨诺夫（D. A.）与阿尼金（V. M.）（2019）。〈*The Sixth Congress of Russian Physicists in Saratov (August 15, 1928)*〉（《1928年8月萨拉托夫第六届俄罗斯物理学家大会》）。《萨拉托夫大学学报·物理学新系列》（*Izvestiya of Saratov University. New Series. Series: Physics*），第19卷（第2期）：153–161。doi:10.18500/1817-3020-2019-19-2-153-161。

\item 辛格（Rajinder Singh）与里斯（Falk Riess）（2001）。〈*The Nobel Prize for Physics in 1930 – A close decision?*〉（《1930年诺贝尔物理学奖——一次接近的决策？》）。《伦敦皇家学会记录与文献》（*Notes and Records of the Royal Society of London*），第55卷（第2期）：267–283。doi:10.1098/rsnr.2001.0143。S2CID 121955580。

\item 法贝林斯基（Immanuil L. Fabelinskiĭ）（2003年10月31日）。〈*The discovery of combination scattering of light in Russia and India*〉（《俄罗斯与印度的光组合散射发现》）。《物理学进展》（*Physics-Uspekhi*），第46卷（第10期）：1105–1112。doi:10.1070/PU2003v046n10ABEH001624。S2CID 250862316。原文存档于2021年8月12日。检索于2020年3月13日。

\item 法贝林斯基（I. L. Fabelinskii）（1990）。〈*Priority and the Raman effect*〉（《优先权与拉曼效应》）。《自然》（*Nature*），第343卷（第6260期）：686。Bibcode:1990Natur.343..686F。doi:10.1038/343686a0。S2CID 4340367。

\item 克里希南（K. S. Krishnan）（1928）。〈*Influence of Temperature on the Raman Effect*〉（《温度对拉曼效应的影响》）。《自然》，第122卷（第3078期）：650。Bibcode:1928Natur.122..650K。doi:10.1038/122650b0。S2CID 4107416。

\item 克里希南（K. S. Krishnan）（1928）。〈*The Raman Effect in Crystals*〉（《晶体中的拉曼效应》）。《自然》，第122卷（第3074期）：477–478。Bibcode:1928Natur.122..477K。doi:10.1038/122477a0。S2CID 4095088。

\item 克里希南（K. S. Krishnan）（1928）。〈*The Raman Effect in X-ray Scattering*〉（《X射线散射中的拉曼效应》）。《自然》，第122卷（第3086期）：961–962。Bibcode:1928Natur.122..961K。doi:10.1038/122961c0。S2CID 4103299。

\item 拉马塞尚（S. Ramaseshan）（1998）。〈*A conversation with K. S. Krishnan on the story of the discovery of the Raman effect*〉（《与K.S.克里希南谈拉曼效应发现的故事》）。《当代科学》（*Current Science*），第75卷（第11期）：1265–1272。JSTOR 24101925。

\item 瓦利（Kameshwar C. Wali）（1991）。《*Chandra: A Biography of S. Chandrasekhar*》（《钱德拉：钱德拉塞卡传记》）。芝加哥：芝加哥大学出版社。ISBN 978-0-226-87054-0。OCLC 21297960。

\item 拉马斯瓦米（Karthik Ramaswamy）（2019年12月12日）。〈*When Raman Brought Born to Bangalore – Connect with IISc*〉（《当拉曼将玻恩带到班加罗尔——与印度科学研究所的联系》）。《Connect》。原文存档于2020年11月7日。检索于2020年3月15日。

\item 玻恩（Max Born）（1965）。〈*Recollections of Max Born II. What I Did as a Physicist*〉（《马克斯·玻恩回忆录 II：作为物理学家的我》）。《原子科学家公报》（*Bulletin of the Atomic Scientists*），第21卷（第8期）：9–13。Bibcode:1965BuAtS..21h...9B。doi:10.1080/00963402.1965.11454843。

\item 卡尔多纳（M. Cardona）与马克斯（W. Marx）（2008年7月1日）。〈*Max Born and his legacy to condensed matter physics*〉（《马克斯·玻恩及其对凝聚态物理的遗产》）。《物理学年鉴》（*Annalen der Physik*），第17卷（第7期）：497–518。Bibcode:2008AnP...520..497C。doi:10.1002/andp.200810304。S2CID 121440005。

\item 苏尔（Abha Sur）（1999）。〈*Aesthetics, Authority, and Control in an Indian Laboratory: The Raman-Born Controversy on Lattice Dynamics*〉（《印度实验室中的美学、权威与控制：拉曼—玻恩关于晶格动力学的争论》）。《伊西斯》（*Isis*），第90卷（第1期）：25–49。doi:10.1086/384240。JSTOR 237473。S2CID 144805021。

\item 达斯古普塔（Deepanwita Dasgupta）（2010年1月11日）。〈*Progress in Science and Science at the Non-Western Peripheries*〉（《科学的进步与非西方边缘地区的科学发展》）。《自发生成：科学史与哲学期刊》（*Spontaneous Generations: A Journal for the History and Philosophy of Science*），第3卷（第1期）：142–157。doi:10.4245/sponge.v3i1.6575。

\item 辛格（Ravinder Singh）。〈*Sir C.V. Raman, Dame Kathleen Lonsdale and their Scientific Controversy due to the Diffuse Spots in X-ray Photographs*〉（《C.V.拉曼爵士与凯瑟琳·朗斯代尔女爵：因X射线照片弥散斑点引发的科学争论》）（PDF）。《印度科学史杂志》（*Indian Journal of History of Science*），第37卷（第3期）：267–290。PDF存档于2016年9月21日。检索于2016年6月30日。

\item 辛格（Rajinder Singh）（2008）。〈*Max Born's Role in the Lattice Dynamic Controversy*〉（《马克斯·玻恩在晶格动力学争论中的角色》）。《半人马座》（*Centaurus*），第43卷（第3–4期）：260–277。doi:10.1111/j.1600-0498.2000.cnt430306.x。

\item 埃鲁南（Earunan）（2018年3月4日）。〈*Jawaharlal Nehru and C. V. Raman: Nehru's vision is more important for Science in India, not Raman's!*〉（《尼赫鲁与C.V.拉曼：对印度科学而言，尼赫鲁的愿景比拉曼更重要！》）。*earunan.com*。原文存档于2020年9月18日。检索于2020年3月15日。

\item 辛格（Rajinder Singh）与里斯（Falk Riess）（2013）。〈*Belated Nobel Prize for Max Born F.R.S.*〉（《迟到的诺贝尔奖：马克斯·玻恩院士》）（PDF）。《印度科学史杂志》，第48卷：79–104。PDF原文存档于2016年7月1日。检索于2016年7月1日。

\item 帕拉梅斯瓦兰（Uma Parameswaran）（2011）。*前引书（Op. cit.）*，第222页。OCLC 772714846。

\item 马尔霍特拉（Inder Malhotra）（2014）。〈*C. V. Raman and the Bharat Ratna*〉（《C.V.拉曼与“印度之光”勋章》）。*freedomfirst.in*。原文存档于2020年10月20日。检索于2020年9月11日。

\item “**Remembering CV Raman's Wit and the Time he Tricked Nehru into Believing Copper is Gold**”（《回忆C.V.拉曼的机智：他曾哄骗尼赫鲁相信铜是金》）。*News18*。2018年11月21日。原文存档于2020年10月27日。检索于2020年3月15日。

\item 斯里坎特（B. R. Srikanth）（2017年2月28日）。〈*No Raman effect: How his dream died a quiet death*〉（《没有“拉曼效应”：他梦想的无声消逝》）。《德干纪事报》（*Deccan Chronicle*）。原文存档于2020年10月23日。检索于2020年3月15日。

\item “**The ups and downs of a science city**”（《科学城的兴衰起伏》）。*downtoearth.org.in*。1994年5月31日。原文存档于2020年10月31日。检索于2020年3月15日。

\item 克里希纳（V. V. Krishna）与哈德里亚（Binod Khadria）（1997）。〈*Phasing Scientific Migration in the Context of Brain Gain and Brain Drain in India*〉（《印度科学迁移的阶段性分析：人才流失与回流背景下的研究》）。《科学、技术与社会》（*Science, Technology and Society*），第2卷（第2期）：347–385。doi:10.1177/097172189700200207。S2CID 143870753。

\item 克里希纳（V. V. Krishna）（2001年6月1日）。〈*Changing policy cultures, phases and trends in science and technology in India*〉（《印度科学与技术政策文化、阶段与趋势的变化》）。《科学与公共政策》（*Science and Public Policy*），第28卷（第3期）：179–194。Bibcode:2001SciPP..28..179K。doi:10.3152/147154301781781525。

\item 巴拉拉姆（P. Balaram）（2009）。〈*Anniversaries at the Academies*〉（《科学院的纪念日》）。《当代科学》（*Current Science*），第96卷（第1期）：5–6。[永久失效链接]

\item 帕拉梅斯瓦兰（Uma Parameswaran）（1999）。*前引书（Op. cit.）*。培生印度教育出版社，第145–147页。ISBN 978-81-317-2818-5。

\item 戈维尔（Girjesh Govil）（2010）。载于达斯古普塔（Uma Dasgupta）（编），《*Science and Modern India: An Institutional History, c. 1784–1947*》（《科学与现代印度：一部制度史（约1784–1947）》）。德里：培生出版社，第143–156页。ISBN 978-93-325-0294-9。OCLC 895913622。

\end{enumerate}